% Generated by scripts/tikz_reviews.py from the figure manifests.
\ReviewFigure{Sampling and spectral overlap}
{figures/1-shannon/sampling-aliasing.png}
{figures/tikz/shannon-sampling-aliasing.pdf}
{The sampling theorem and Poisson formula identify sampling with periodic replication of the spectrum.}
{Unit sample spacing is used. These are exact compactly supported cubic B-spline spectra of sinc(x/a)\textasciicircum{}4, with a=4 and a=2. Replicas are spaced by 2 pi; the shaded Nyquist interval is [-pi, pi]. Ideal reconstruction retains the sum inside that interval. The right-hand result differs from the original spectrum.}
{shannon-sampling-aliasing}
\ReviewFigure{The cardinal sine}
{figures/1-shannon/sinc.png}
{figures/tikz/shannon-sinc.pdf}
{The reconstruction kernel in Shannon interpolation uses the normalized cardinal sine.}
{The value at zero is the continuous extension, equal to one. Every nonzero integer is a zero. The plotted normalization is sin(pi x)/(pi x), as in the theorem.}
{shannon-sinc}
\ReviewFigure{Cardinal B-splines of degrees zero, one, and two}
{figures/1-shannon/spline.png}
{figures/tikz/shannon-splines.pdf}
{The chapter defines phi zero as the centered unit box and obtains each subsequent spline by convolution with that box.}
{Supports and heights follow the convolution definition exactly. The quadratic spline peaks at 3/4, so its integer translates are not an interpolating basis with the samples as coefficients. Endpoint values of the box follow the closed-interval convention in the text.}
{shannon-splines}
\ReviewFigure{Two frequencies with identical samples}
{figures/1-shannon/aliasing.png}
{figures/tikz/shannon-aliasing.pdf}
{A periodic sinusoidal example explains aliasing through discrete spectral lines.}
{For unit sample spacing, cos(pi x/2) and cos(3 pi x/2) agree at every integer. Frequencies are identified modulo 2 pi; the high-frequency pair folds to the low-frequency pair. This example is periodic and is not an L2 signal on the real line.}
{shannon-aliasing}
\ReviewFigure{Uniform scalar quantization}
{figures/1-shannon/quantizer.pdf}
{figures/tikz/shannon-quantizer.pdf}
{The quantizer returns an integer index; dequantization multiplies that index by the step size T.}
{The horizontal axis is u/T and the vertical axis is the integer v. Filled left endpoints and open right endpoints implement v-1/2 <= u/T < v+1/2 exactly. The diagonal is the unquantized value; the discrepancy is at most one half in these normalized coordinates.}
{shannon-quantizer}
